\usepackage{iftex}
\ifXeTeX
  \usepackage{fontspec}
\else
  \usepackage[utf8]{inputenc}
  \usepackage[T1]{fontenc}
  \usepackage{lmodern}
\fi
\usepackage[spanish,es-tabla,es-nodecimaldot,es-noquoting,es-noshorthands]{babel}
\usepackage[a4paper,margin=2.5cm]{geometry}
\usepackage{microtype}
\usepackage[htt]{hyphenat}
\usepackage{graphicx}
\usepackage{float}
\usepackage[section]{placeins}
\usepackage{booktabs,tabularx,array,longtable,calc}
\usepackage{enumitem}
\usepackage[dvipsnames]{xcolor}
\usepackage[font=small,labelfont=bf]{caption}
\captionsetup[figure]{labelformat=empty}
\usepackage{listings}
\usepackage{url}
\usepackage[hidelinks]{hyperref}
\hypersetup{pdftitle={Práctica 2. Guion de prácticas CLIPS},pdfauthor={Pablo Pardo Cotos},pdfsubject={Intelixencia Artificial}}
\graphicspath{{img/}}
\setlength{\parindent}{0pt}
\setlength{\parskip}{4pt}
\setlength{\emergencystretch}{2em}
\setlist{itemsep=1pt,topsep=2pt,leftmargin=*}
\renewcommand{\arraystretch}{1.15}
\newcolumntype{Y}{>{\raggedright\arraybackslash}X}
\setcounter{tocdepth}{2}
% Las subsecciones llevan su propio número de ejercicio en el título
\setcounter{secnumdepth}{1}
% Colocación de figuras
\renewcommand{\topfraction}{0.9}
\renewcommand{\bottomfraction}{0.8}
\renewcommand{\textfraction}{0.07}
\renewcommand{\floatpagefraction}{0.75}
\pagestyle{plain}
\urlstyle{same}
% Código CLIPS
\definecolor{azul}{HTML}{17486C}
\definecolor{verde}{HTML}{22685E}
\definecolor{papel}{HTML}{F1F6F9}
\lstdefinelanguage{CLIPS}{
  morekeywords={defrule,deffacts,deffunction,deftemplate,assert,retract,printout,read,
    bind,if,then,else,not,and,or,test,declare,salience,crlf,reset,run,clear,load,batch,facts,agenda,rules},
  sensitive=true,
  morecomment=[l]{;},
  morestring=[b]",
  alsoletter={-?},
}
\lstset{language=CLIPS,basicstyle=\ttfamily\small,keywordstyle=\color{azul}\bfseries,
  commentstyle=\color{verde},stringstyle=\color{Maroon},showstringspaces=false,
  breaklines=true,columns=fullflexible,keepspaces=true,frame=single,
  rulecolor=\color{azul!30},backgroundcolor=\color{papel},aboveskip=6pt,belowskip=6pt,tabsize=3,
  literate={á}{{\'a}}1 {é}{{\'e}}1 {í}{{\'i}}1 {ó}{{\'o}}1 {ú}{{\'u}}1 {ñ}{{\~n}}1 {¿}{{?`}}1}
% Macros que usa el texto convertido desde Markdown
\providecommand{\tightlist}{\setlength{\itemsep}{0pt}\setlength{\parskip}{0pt}}
\makeatletter\@ifundefined{c@none}{\newcounter{none}}{}\makeatother
